\documentclass[11pt]{article}

\title{AI PenTesting Suite}
\usepackage{graphicx}
\usepackage[utf8]{inputenc}
\usepackage{hyperref}
\usepackage{float}
\usepackage{bookmark}
\usepackage[left=2.5cm, right=2.5cm, top=2.5cm, bottom=2.5cm]{geometry}
\usepackage{titlesec}
\usepackage[backend=biber]{biblatex}
\usepackage{tocloft} % Optional: to customize TOC appearance

% Redefine paragraph to behave like a subsection for TOC purposes
\titleformat{\paragraph}
  [runin] % Keep the paragraph style without line break
  {\normalfont\normalsize\bfseries} % Style
  {\theparagraph} % Numbering style
  {1em}
  {}

\titlespacing*{\paragraph}{0pt}{3.25ex plus 1ex minus .2ex}{1em}

\setcounter{secnumdepth}{4}
\setcounter{tocdepth}{4}
\author{Alessandro Luppino}

\addbibresource{references.bib}

\begin{document}
\begin{titlepage}
   \begin{center}
       \vspace*{1cm}
       \vspace{0.5cm}
        {\huge Analisi e documentazione del progetto}
       \begin{figure}[H]
    		\centering
    		\includegraphics[width=0.8\textwidth]{../Resources/Icon/AIPS_Icon2B.png}
    \end{figure}
       \vspace{0.5cm}

       \textbf{Alessandro Luppino}

       \vfill
            
       Politecnico di Bari\\
       Ingegneria del Software
            
       \vspace{0.1cm}
   \end{center}
\end{titlepage}
\tableofcontents
\newpage
\section{Introduzione}
Il progetto consiste nello sviluppo di una chat da terminale per sistemi Linux, 
che consente di interagire con un modello linguistico avanzato (LLM)\autocite{YAO2024100211}, 
appositamente integrato con strumenti di sicurezza esistenti.\\
L'obiettivo principale è l'applicazione del modello di intelligenza artificiale 
nell'ambito della cybersecurity, con lo scopo di assistere l'utente nella prima 
fase di penetration testing\autocite{app13126986} fornendo un'esperienza interattiva.\\
\section{Contesto e stato dell'arte}
Nel panorama attuale della cybersecurity, l’intelligenza artificiale ha assunto un
ruolo sempre più centrale\autocite{Princess_Egbuna_2021}\autocite{zhang2024llmsmeetcybersecuritysystematic}, con modelli avanzati che supportano analisi delle minacce, 
rilevamento delle vulnerabilità e simulazioni di attacchi\autocite{GetPwnd2023}. Molte soluzioni disponibili 
sul mercato fanno uso di modelli altamente specializzati, addestrati su dataset specifici 
relativi alla sicurezza informatica, che garantiscono performance ottimizzate. 
Tuttavia, questi strumenti presentano spesso delle limitazioni, come la necessità 
di accedere a servizi remoti, restrizioni nel numero di richieste effettuabili e costi 
associati all’uso di modelli complessi.\\
\textbf{AI PenTesting Suite} si inserisce in questo contesto differenziandosi 
dai tradizionali servizi AI in quanto prevede l'hosting del modello sulla macchina dell'utente 
(di consueto si utilizza un'architettura Client-Server), e l'integrazione con strumenti di cybersecurity affermati.
\section{Requisiti}
\subsection{Funzionali}
\begin{itemize}
    \item \textbf{Interfaccia conversazionale: }il sistema deve consentire agli utenti di interagire tramite una chat da terminale.
    \item \textbf{Scansione delle vulnerabilità: }deve eseguire scansioni delle infrastrutture IT utilizzando strumenti come Nmap e analizzarne i risultati per identificare vulnerabilità, assegnando un livello di rischio.
    \item \textbf{Ricerca delle vulnerabilità: }deve consentire agli utenti di cercare informazioni su vulnerabilità specifiche, restituendo dettagli e suggerimenti.
\end{itemize}
\subsection{Non Funzionali}
\begin{itemize}
    \item \textbf{Prestazioni: }il sistema deve elaborare i dati di una scansione e generare risposte in tempi ragionevoli.
    \item \textbf{Scalabilità: }deve essere facilmente estendibile per includere nuove integrazioni.
    \item \textbf{Usabilità: }l’output deve essere chiaro e diretto.
    \item \textbf{Affidabilità: }il sistema deve fornire risultati accurati relativi alle vulnerabilità rilevate.
\end{itemize}
\section{Modello di processo}
Il progetto è stato sviluppato seguendo un approccio iterativo e incrementale, 
integrato con una chiara definizione dei requisiti funzionali e non funzionali. 
Questo approccio si è rivelato efficace per affrontare lo sviluppo di un sistema complesso, 
che è stato suddiviso in moduli per facilitare il testing durante lo sviluppo e gli aggiornamenti futuri. 
L’iterazione costante ha permesso di soddisfare i requisiti specificati in modo graduale, 
garantendo un miglioramento continuo e un risultato stabile.
\begin{figure}[H]
    \centering
    \includegraphics[width=1\textwidth]{../Resources/Diagrams/Gantt.png}
    Il Gantt iniziale del progetto.
\end{figure}
\section{Analisi dei Costi} \label{AdC}
I risultati di questa analisi sono quelli che hanno determinato le scelte che verranno poi discusse.
\subsection{Hosting del modello} \label{HdM}
L’hosting di un modello da 1.5B parametri su una piattaforma come \href{https://huggingface.co/spaces}{Hugging Face Spaces} 
richiede risorse significative per garantire prestazioni accettabili durante le interazioni: 
almeno una GPU con 16-24 GB di memoria per essere caricato e mantenere tempi di risposta adeguati.\\
\href{https://huggingface.co/spaces}{Hugging Face Spaces} offre un piano gratuito, ma con risorse limitate, 
non adatte a modelli di questa scala. Questo è stato verificato durante i test, 
dove il sistema risultava così lento da essere inutilizzabile.\\\\
Per accedere a risorse adeguate, è quindi necessario un piano premium o l’uso di servizi cloud 
come \href{https://aws.amazon.com/it/}{AWS}, \href{https://cloud.google.com/}{Google Cloud} o 
\href{https://portal.azure.com/}{Azure}.\\
su \href{https://huggingface.co/spaces}{Hugging Face Spaces}, un'istanza GPU T4 dedicata costa circa \$0.60-\$1.00 all'ora, 
mentre un'istanza con una GPU A100 (più adatta per i carichi richiesti dal 
progetto) può salire a \$2.00-\$4.00 all'ora.\\
Supponendo un utilizzo continuo di 24 ore al giorno in modo che il modello sia sempre disponibile, 
i costi variano da \$432 a \$2880 al mese a seconda delle risorse selezionate.\\
\textit{NOTA: Ci sono alternative più economiche di quella analizzata,
piattaforme come Google Colab Pro+ o Paperspace offrono GPU a tariffe ridotte, 
ma la scalabilità e la disponibilità per un’applicazione a lungo termine restano 
problematiche e l'hosting su uno di questi servizi è troppo dispendioso.}
\newpage
\subsection{Fine-tuning del modello}
\subsubsection{In locale}
Il fine-tuning di un modello da 1.5B parametri è un processo computazionalmente intenso
che richiede hardware e risorse software specifiche\autocite{j2024finetuningllmenterprise}:
\begin{itemize}
    \item \textbf{GPU con almeno 24 GB di VRAM:} quantità minima necessaria per 
    gestire il training in batch ragionevoliS.
    \item \textbf{CPU potente:} per gestire il pre-processing dei dati e il training in parallelo.
    \item \textbf{Ampio spazio di archiviazione:} per caricare dataset e checkpoint intermedi.
\end{itemize}
La mia macchina non soddisfa nessuno di questi requisiti, e l’acquisto di hardware dedicato per il fine-tuning comporterebbe spese troppo elevate.\\
\textit{NOTA: Anche se fosse possibile suddividere i batch per adattarsi a una GPU con meno memoria, 
i tempi di elaborazione crescerebbero esponenzialmente, rendendo il processo impraticabile.}
\subsubsection{In Cloud}
Un'alternativa più economica per il fine-tuning sono le stesse piattaforme Cloud discusse nell'hosting del modello \ref{HdM}.
Il training completo di un modello da 1.5B parametri può richiedere sulle 48-72 ore di calcolo 
su una GPU V100. Un’istanza GPU V100 su \href{https://aws.amazon.com/it/}{AWS} costa circa \$2.50-\$3.00 all’ora, 
mentre una GPU A100 può costare fino a \$4.00-\$5.00 all’ora.
Il costo totale varia quindi da \$120 a \$360 per il fine-tuning 
(escludendo costi per archiviazione e trasferimento dati).\\
\textit{NOTA: il fine-tuning richiede anche dataset specifici e pre-elaborati, il che aggiunge tempo e possibili costi.}
\section{Problemi riscontrati e Scelte tecniche}
Uno dei principali problemi affrontati durante lo sviluppo del progetto, è stata 
l’impossibilità tecnica di effettuare il fine-tuning del modello AI su dati 
specifici del dominio della cybersecurity\autocite{dataset}. Questo limite, 
legato principalmente a risorse computazionali insufficienti e a costi elevati \ref{AdC}, 
ha portato a scegliere un modello general-purpose già pre-addestrato, 
accettando compromessi in termini di specificità delle risposte.\\
Per mitigare gli effetti negativi di questa scelta e le "allucinazioni" tipiche di qualsiasi
LLM, il sistema è stato, attualmente, integrato con:
\begin{itemize}
    \item \textbf{nmap\autocite{nmap_docs}}: per effettuare scansioni di reti e porte, fornendo dettagli utili per valutazioni di sicurezza.
    \item \textbf{nvdilb\autocite{nvdlib_docs}}: una libreria python per interagire con il \href{https://nvd.nist.gov/}{NVD (National Vulnerability Database)}, permettendo la ricerca e il controllo di eventuali CVE (Common Vulnerabilities and Exposures) associate ai servizi.
\end{itemize}
E prevede l'integrazione di ulteriori strumenti in futuro. (vedi Sviluppi Futuri)\\
Queste integrazioni permette di compensare la mancanza di conoscenza specifica del modello, 
migliorando l'accuratezza, l'affidabilità e la ripetibilità delle risposte.\\
Un altro aspetto distintivo del progetto è il prima citato hosting locale sulla macchina dell’utente, 
che garantisce piena autonomia e libertà di utilizzo del sistema senza restrizioni. 
Tuttavia, questa scelta comporta un aumento del carico computazionale sulla macchina stessa, 
richiedendo un compromesso tra performance del sistema e risorse disponibili. 
La decisione di privilegiare un’architettura locale, oltre a fornire una soluzione completamente 
gratuita, riflette l’obiettivo di creare uno strumento flessibile, 
senza dipendenze da infrastrutture esterne, adatto a scenari dove 
la privacy e la disponibilità del modello giocano un ruolo critico.\\
\textit{Ad Es. è possibile conversare con il modello anche in assenza di connessione internet seppure 
si svanno a sacrificare alcune funzionalità.}
\section{Prototipo}
Il primo prototipo del progetto è stato presentato in data 29/11/2024.\\ 
In questa fase, il sistema implementava le funzionalità di base, soddisfacendo parzialmente 
i requisiti funzionali e non funzionali precedentemente definiti. Il prototipo era già operativo
e in grado di svolgere alcune funzionalità più complesse come la scansione di reti tramite Nmap.
Il problema principale riscontrato è stato il tempo di risposta del modello, soprattutto durante le scansioni
i tempi di attesa superavano il minuto, rendendo l'esperienza utente non ottimale.\\
\begin{figure}[H]
    \centering
    \includegraphics[width=1\textwidth]{../Resources/Diagrams/CasiDUso.jpg}
    Il diagramma dei casi d'uso seguito anche dal prototipo.
\end{figure}
\section{Stato Attuale}
Dopo aver analizzato attentamente i limiti del prototipo, sono state effettuate numerose migliorie
all'esperienza utente, tra cui un'interfaccia grafica molto più pulita che permette
al modello di formattare le risposte, l'aggiunta di un sistema di cronologia in modo che l'utente
possa automaticamente ripetere comandi passati più facilmente, e la completa rivisitazione del
sistema che gestiva l'input e l'output del modello, per permette di vedere in tempo reale i messaggi
del modello, riducendo notevolmente i tempi di attesa.\\
Oltre a questo sono state aggiunte nuove funzionalità come la ricerca di CVE tramite NVD e resa più
efficace l'integrazione con nmap.\\
\begin{figure}[H]
    \centering
    \includegraphics[width=1\textwidth]{../Resources/Diagrams/Architettura.jpg}
\end{figure}
\begin{figure}[H]
    \centering
    \includegraphics[width=1\textwidth]{../Resources/Diagrams/Sequenze.jpg}
    Architettura e diagramma delle sequenze della versione finale del progetto.
\end{figure}

\section{Sviluppi Futuri}
Gli sviluppi futuri del progetto mirano a potenziarne ulteriormente le funzionalità e 
migliorare l’esperienza utente.
È previsto l’aggiungimento di nuovi moduli, come un Modulo 
di Attacco, che suggerirà exploit e tecniche di attacco specifiche per le vulnerabilità 
rilevate, e un Modulo di Report, capace di generare report dettagliati e professionali 
utili per la documentazione del lavoro svolto.\\
Inoltre, si intende implementare un sistema per salvare le conversazioni e consentire agli 
utenti di riprenderle in futuro, garantendo continuità nelle analisi.\\
Sul fronte dell’usabilità, lo sviluppo di un’interfaccia grafica più intuitiva rappresenta 
una priorità, insieme a funzionalità che migliorino l’interazione, come la possibilità di 
interrompere il modello durante la generazione di output non soddisfacenti.\\
Infine, se i fondi lo permetteranno, si prevede di effettuare il fine-tuning del modello AI. 
Questo consentirebbe di adattare il sistema a scenari specifici della cybersecurity, 
migliorandone la precisione e riducendo ulteriormente il rischio di allucinazioni.
\\
\\
\section*{Stack Tecnico}

Il modello utilizzato per il progetto è \href{https://huggingface.co/Qwen/Qwen2.5-1.5B-Instruct}{\texttt{Qwen2.5-1.5B-Instruct}}, un modello AI general-purpose con 1.5 miliardi di parametri, scelto per le sue capacità avanzate di generazione e comprensione del linguaggio naturale. Lo sviluppo è stato effettuato in Python, versione \texttt{3.12.8}, con il supporto delle seguenti librerie:

\begin{itemize}
    \item \texttt{json}, \texttt{time}, \texttt{os}, \texttt{re}, \texttt{socket}: per la gestione di operazioni di base del sistema e manipolazione dei dati.
    \item \texttt{nvdlib}, \texttt{nmap}, \texttt{whois}: per l'integrazione con strumenti e librerie di cybersecurity.
    \item \texttt{prompt\_toolkit} (moduli: \texttt{PromptSession}, \texttt{HTML}, \texttt{KeyBindings}, \texttt{FileHistory}, \texttt{run\_in\_terminal}, \texttt{AutoSuggestFromHistory}): per implementare l'interfaccia conversazionale interattiva.
    \item \texttt{rich} (moduli: \texttt{Console}, \texttt{Live}, \texttt{Markdown}, \texttt{Spinner}): per fornire un output visuale avanzato e leggibile.
    \item \texttt{threading.Thread}: per gestire l'esecuzione di task concorrenti.
    \item \texttt{transformers} (moduli: \texttt{AutoTokenizer}, \texttt{AutoModelForCausalLM}, \texttt{TextIteratorStreamer}): per il caricamento e l'interazione con il modello AI.
    \item \texttt{torch} (\textit{PyTorch}): come backend per il machine learning, necessario per il funzionamento del modello.
\end{itemize}

Test e sviluppo sono stati eseguiti su una macchina con \texttt{AMD Ryzen 5 7535U} e \texttt{16 GB} di RAM.

\newpage
\printbibliography
\end{document}